%==============================================================================
% CHAPTER 5: What Is... an Equivariant Neural Network?
%==============================================================================
\chapter{What Is... an Equivariant Neural Network?}
\label{ch:equivariant} \cite{lim2023equivariant}

\begin{flushright}
\textit{Lek-Heng Lim, Bradley J. Nelson}\\
Communicated by \textit{Notices} Associate Editor Emilie Purvine
\end{flushright}

\epigraph{``An equivariant neural network is an equivariant map constructed from alternately composing equivariant linear maps with nonlinear ones.''}{--- Lim \& Nelson} \cite{lim2023equivariant}

\begin{abstract}
Symmetry is everywhere in science: molecular structures, physical laws, images. Standard neural networks ignore symmetry, \index{symmetry in neural networks}wasting capacity learning what physics already dictates. Equivariant neural networks bake symmetry into the architecture: if the input transforms, the output transforms predictably. This chapter gives the mathematical foundation: group representation\index{group representations}s, intertwining operators, and the universal approximation theorem for equivariant networks, \cite{lim2023equivariant}.
\end{abstract}

%==============================================================================
\section{Introduction: Why Symmetry Matters}
\label{sec:eq:intro}

A molecule rotated in space is the \emph{same} molecule. A cat translated in an image is the same cat. The laws of physics are translation/rotation invariant. Yet a standard CNN must \emph{learn} translation invariance from data---a waste of parameters and samples.

An \textbf{equivariant neural network} guarantees: \cite{lim2023equivariant}
\[
f(g \cdot x) = g \cdot f(x) \quad \forall g \in G,
\]
where $G$ is a symmetry group acting on input space $\mathcal{X}$ and output space $\mathcal{Y}$, \cite{lim2023equivariant}.

\begin{itemize}
    \item \textbf{Invariance}: $f(g \cdot x) = f(x)$ (e.g., classification: rotated cat $\to$ ``cat'')
    \item \textbf{Equivariance}: $f(g \cdot x) = g \cdot f(x)$ (e.g., segmentation: rotated input $\to$ rotated mask)
    \item \textbf{Covariance}: Output transforms by a \emph{different} representation of $G$.
\end{itemize}

%==============================================================================
\section{Group Representations}
\label{sec:eq:reps}

\begin{definition}[Group Representation]
A \textbf{representation} of a group $G$ on a vector space $V$ is a homomorphism $\rho: G \to \mathrm{GL}(V)$. For $g \in G$, $\rho(g)$ is an invertible linear map on $V$, \cite{lim2023equivariant}.
\end{definition}

Key examples for deep learning:
\begin{itemize}
    \item \textbf{Translation} ($G = \mathbb{Z}^2$ or $\mathbb{R}^2$): $\rho(g)_{ij} = x_{i-g}$ (shift)
    \item \textbf{Rotation} ($G = \mathrm{SO}(2)$ or $\mathrm{SO}(3)$): $\rho(g) = R_g$ (rotation matrix)
    \item \textbf{Permutation} ($G = S_n$): $\rho(\sigma)_{ij} = \delta_{i, \sigma(j)}$
    \item \textbf{Permutation + Sign} ($G = S_n \times \{\pm 1\}^n$): for fermionic wavefunctions
\end{itemize}

\begin{theorem}[Schur's Lemma]
Let $\rho_1: G \to \mathrm{GL}(V_1)$, $\rho_2: G \to \mathrm{GL}(V_2)$ be irreducible representations. Any linear map $T: V_1 \to V_2$ satisfying $T \rho_1(g) = \rho_2(g) T$ for all $g$ is either zero or an isomorphism. If $V_1 = V_2$ and $\rho_1 = \rho_2$ irreducible, then $T = \lambda I$.
\end{theorem}

\textbf{Consequence}: Equivariant linear layers are highly constrained! For irreps, they are just scalar multiples of the identity.

%==============================================================================
\section{Equivariant Linear Layers}
\label{sec:eq:linear}

We want linear $\bm{W}: V_1 \to V_2$ such that $\bm{W} \rho_1(g) = \rho_2(g) \bm{W}$.

\begin{example}[Convolution = Translation Equivariance]
Let $G = \mathbb{Z}^2$ act on images by translation. An equivariant linear map is exactly a \textbf{convolution}: \cite{lim2023equivariant}
\[
(\bm{W} * x)[i] = \sum_j w[i-j] x[j].
\]
Proof: $(T_g x)[i] = x[i-g]$. Then $(T_g \bm{W} * x)[i] = \sum_j w[i-j] x[j-g] = (\bm{W} * T_g x)[i]$.
\end{example}

\begin{theorem}[General Solution for Linear Equivariant Maps]
For finite $G$, the space of equivariant linear maps between representations $\rho_1, \rho_2$ is \cite{lim2023equivariant}
\[
\mathrm{Hom}_G(\rho_1, \rho_2) = \left\{ \frac{1}{|G|} \sum_{g \in G} \rho_2(g) M \rho_1(g^{-1}) : M \in \mathbb{R}^{d_2 \times d_1} \right\}.
\]
This is the \emph{projection} onto the equivariant subspace (Reynolds operator), \cite{lim2023equivariant}.
\end{theorem}

%==============================================================================
\section{Equivariant Nonlinearities}
\label{sec:eq:nonlinear}

Linear equivariance \index{equivariant networks}is well understood. Nonlinearities break equivariance unless carefully chosen.

\begin{proposition}
A pointwise nonlinearity $\sigma: \mathbb{R} \to \mathbb{R}$ is $G$-equivariant iff it commutes with the action of $G$ on $\mathbb{R}$, \cite{lim2023equivariant}.
\end{proposition}

For permutation groups $S_n$ acting by permuting coordinates, \textbf{any} pointwise $\sigma$ is equivariant (since $\sigma(x_{\pi(i)}) = \sigma(x)_{\pi(i)}$). This is why ReLU works in standard MLPs and CNNs! \cite{lim2023equivariant}

For rotation groups $\mathrm{SO}(3)$, pointwise $\sigma$ on vector components is \emph{not} equivariant (rotation mixes components). Solutions: \cite{lim2023equivariant}
\begin{enumerate}
    \item \textbf{Scalarization}: Apply $\sigma$ to \emph{invariants} (norms, dot products)
    \item \textbf{Gated nonlinearities}: $\sigma(\bm{x}) = \sigma_{\text{scalar}}(\|\bm{x}\|) \cdot \bm{x}/\|\bm{x}\|$
    \item \textbf{Tensor product}: Work in irreps; Clebsch-Gordan coefficients give equivariant polynomial maps \cite{lim2023equivariant}
\end{enumerate}

%==============================================================================
\section{Architecture: Equivariant Neural Networks}
\label{sec:eq:arch}

\begin{definition}[Equivariant Neural Network]
A feedforward network where:
\begin{enumerate}
    \item Each layer $\bm{x}^{(\ell+1)} = \sigma(\bm{W}^{(\ell)} \bm{x}^{(\ell)} + \bm{b}^{(\ell)})$
    \item $\bm{W}^{(\ell)}$ is $G$-equivariant (intertwiner) \cite{lim2023equivariant}
    \item $\sigma$ is $G$-equivariant (pointwise on scalars, gated on vectors) \cite{lim2023equivariant}
    \item $\bm{b}^{(\ell)}$ is $G$-invariant (typically zero or learned invariant features)
\end{enumerate}
\end{definition}

The composition of equivariant maps is equivariant, so the whole network is equivariant, \cite{lim2023equivariant}.

\begin{example}[$\mathrm{SO}(3)$-Equivariant Network for Molecules]
\begin{enumerate}
    \item Input: atomic positions $\bm{r}_i \in \mathbb{R}^3$, types $z_i$
    \item Embed types as scalars (invariant features)
    \item Spherical harmonics $Y_l^m(\bm{r}_i - \bm{r}_j)$ as equivariant features (transform as $D^l(R)$) \cite{lim2023equivariant}
    \item Tensor product layers using Clebsch-Gordan coefficients
    \item Gated nonlinearities preserve equivariance
    \item Output: scalar (energy), vector (forces), tensor (stress)
\end{enumerate}
This is the architecture behind \textbf{NequiP}, \textbf{MACE}, \textbf{SE(3)\index{SE(3) equivariance}-Transformer}.
\end{example}

%==============================================================================
\section{Code: Implementing an Equivariant Layer}
\label{sec:eq:code}

\begin{lstlisting}[style=python,caption=SO(2)-equivariant convolution using steerable filters] \cite{lim2023equivariant}
import torch
import torch.nn as nn
import torch.nn.functional as F
from math import factorial

class SO2Conv(nn.Module):
    """
    Steerable CNN for SO(2) equivariance.
    Features transform as Fourier modes: f_l(r, theta) = sum_m f_{l,m}(r) e^{im theta}
    """
    def __init__(self, in_channels, out_channels, max_freq, kernel_size=3):
        super().__init__()
        self.max_freq = max_freq
        self.in_channels = in_channels
        self.out_channels = out_channels
        
        # Learnable radial profiles for each frequency
        # For each input freq l_in and output freq l_out, we need a radial kernel
        # l_out = l_in + m where m is the filter's frequency
        self.radial_mlps = nn.ModuleDict()
        
        for l_in in range(max_freq + 1):
            for l_out in range(max_freq + 1):
                m = l_out - l_in  # filter frequency
                if abs(m) <= max_freq:
                    key = f"{l_in}_{l_out}"
                    # Radial profile: small MLP taking radius as input
                    self.radial_mlps[key] = nn.Sequential(
                        nn.Linear(1, 16), nn.ReLU(),
                        nn.Linear(16, in_channels * out_channels)
                    )
    
    def forward(self, x):
        # x: [B, C_in, L_max+1, 2, H, W] 
        #   complex features: real/imag for each frequency
        B, C_in, L, _, H, W = x.shape
        
        # Convert to polar coordinates for convolution
        # This is a simplified version; real implementation uses FFT
        out = torch.zeros(B, self.out_channels, self.max_freq+1, 2, H, W, device=x.device)
        
        for l_in in range(self.max_freq + 1):
            for l_out in range(self.max_freq + 1):
                key = f"{l_in}_{l_out}"
                if key not in self.radial_mlps:
                    continue
                
                # Get radial weights
                # In practice, we evaluate radial profile on grid and do convolution via FFT
                pass  # Full implementation requires careful FFT handling
        
        return out

# Simpler: Scalar + Vector fields for SO(3)
class SO3Conv(nn.Module):
    """Tensor field network style: scalars (l=0) and vectors (l=1)"""
    def __init__(self, in_scalar, in_vector, out_scalar, out_vector):
        super().__init__()
        # Scalar -> Scalar: standard conv
        self.ss = nn.Conv2d(in_scalar, out_scalar, 3, padding=1)
        # Scalar -> Vector: multiply by learned vector field
        self.sv = nn.Conv2d(in_scalar, out_vector * 3, 3, padding=1)
        # Vector -> Scalar: dot product with learned vector
        self.vs = nn.Conv2d(in_vector * 3, out_scalar, 3, padding=1)
        # Vector -> Vector: linear combination + cross product
        self.vv = nn.Conv2d(in_vector * 3, out_vector * 3, 3, padding=1)
    
    def forward(self, scalars, vectors):
        # scalars: [B, C_s, H, W]
        # vectors: [B, C_v, 3, H, W]  (3 components)
        B, C_v, _, H, W = vectors.shape
        v_flat = vectors.view(B, C_v * 3, H, W)
        
        out_s = self.ss(scalars) + self.vs(v_flat)
        out_v_flat = self.sv(scalars) + self.vv(v_flat)
        out_v = out_v_flat.view(B, -1, 3, H, W)
        
        return out_s, out_v
\end{lstlisting}

%==============================================================================
\section{Universal Approximation}
\label{sec:eq:univ}

\begin{theorem}[Equivariant Universal Approximation]
Let $G$ be a compact group. An equivariant neural network with: \cite{lim2023equivariant}
\begin{itemize}
    \item Equivariant linear layers spanning $\mathrm{Hom}_G(\rho_1, \rho_2)$
    \item Equivariant nonlinearities separating orbits
\end{itemize}
can approximate any continuous $G$-equivariant function $f: V_1 \to V_2$ uniformly on compact sets, \cite{lim2023equivariant}.
\end{theorem}

\textbf{Key insight}: The nonlinearity must \emph{not} be $G$-invariant (e.g., ReLU on scalars works for $S_n$ but not $\mathrm{SO}(3)$). Gated nonlinearities provide the needed separation.

%==============================================================================
\section{Bridge to Chapter 6}
\label{sec:eq:bridge}

Equivariant networks constrain the \textbf{symmetry} of the affine spline tessellation (Chapter 4). Topological deep learning (Chapter 6) studies the \textbf{shape} of the decision boundaries and feature spaces using algebraic topology invariants. Both are structural priors: symmetry reduces parameters; topology measures complexity, \cite{lim2023equivariant}.

%==============================================================================
\section{Further Reading}
\label{sec:eq:refs}

\begin{itemize}
    \item Lim, Nelson (2023). \textit{What Is... an Equivariant Neural Network?} Notices, 70(4).
    \item Cohen, Welling (2016). \textit{Group Equivariant CNNs}. ICML.
    \item Thomas et al. (2018). \textit{Tensor Field Networks}. NeurIPS.
    \item Fuchs et al. (2020). \textit{SE(3)-Transformers}. ICML.
    \item Geiger, Smidt (2022). \textit{e3nn: Euclidean Neural Networks}. \url{https://github.com/e3nn/e3nn}
\end{itemize}

%==============================================================================
\section{Exercises}
\label{sec:eq:exercises}

\begin{exercise}
Show that a standard CNN with ReLU is translation-equivariant but not rotation-equivariant, \cite{lim2023equivariant}.
\end{exercise}

\begin{exercise}
Implement the Reynolds operator for $G = D_4$ (dihedral group of square) acting on $2 \times 2$ images. Verify it projects any matrix onto the $D_4$-equivariant subspace, \cite{lim2023equivariant}.
\end{exercise}

\begin{exercise}
For $\mathrm{SO}(2)$ acting on $\mathbb{R}^2$ by rotation, find all equivariant linear maps $\mathbb{R}^2 \to \mathbb{R}^2$. What about $\mathbb{R}^2 \to \mathbb{R}$? \cite{lim2023equivariant}
\end{exercise}

\begin{exercise}[Research]
Design an equivariant architecture for the group $G = \mathrm{SE}(3)$ (rigid motions) acting on point clouds. Compare to PointNet++ on ModelNet40, \cite{lim2023equivariant}.
\end{exercise}